\input{sol_pre.tex}
\input{q_2024-10-12_2.tex}
\begin{document}
\input{sol_head.tex}

\Savol{1}
$a+2b-1$ eng katta — $b$ eng katta: $b=24$, $a=1$: \ $1+48-1=48$.
\Javob{D}

\Savol{2}
$81^{0{,}75}=27$, $32^{-\frac25}=\frac14$, $27^{0{,}(3)}=27^{\frac13}=3$, $16^{-0{,}5}=\frac14$, $256^{\frac12}=16$:
\[\frac{27}4-\frac34+16=6+16=22.\]
\Javob{B}

\Savol{3}
$10$ daqiqa $=\frac16$ soat: \ $\dfrac S{120}-\dfrac S{150}=\dfrac16$ $\Rightarrow$ $\dfrac S{600}=\dfrac16$, \ $S=100$ km.
\Javob{A}

\Savol{4}
$1{,}25\cdot1{,}2=1{,}5$ — $50\%$ ga ortgan.
\Javob{B}

\Savol{5}
\[\frac{7^{2m}\cdot11^{2m}}{7^{2m-1}\cdot11^{2m+1}}=\frac7{11}.\]
\Javob{B}

\Savol{6}
$\sqrt2\left(\sqrt6+\sqrt3-\sqrt2-1\right)=2\sqrt3+\sqrt6-2-\sqrt2$ — surat. \ Kasr $\sqrt2$, kvadrati $2$.
\Javob{A}

\Savol{7}
$t=\sqrt[3]2$, $t^3=2$, $\sqrt[3]4=t^2$:
\[\frac3{2t}\left(\frac1{1+t}-\frac1{1-t}\right)=\frac3{2t}\cdot\frac{-2t}{1-t^2}=\frac3{t^2-1}=\frac{3\left(t^4+t^2+1\right)}{t^6-1}=t^4+t^2+1.\]
$t^4+t^2+1+t^2=t^4+2t^2+1=\left(t^2+1\right)^2=\left(\sqrt[3]4+1\right)^2$.
\Javob{C}

\Savol{8}
\[\frac{(x-4)(x-2)}{x-1}-\frac{x-4}{(x-1)(x-2)}=\frac{(x-4)\left((x-2)^2-1\right)}{(x-1)(x-2)}=\frac{(x-4)(x-3)}{x-2}\le0,\quad x\ne1.\]
$x\in(-\infty;1)\cup(1;2)\cup[3;4]$. Natural sonlar: $3$, $4$ — köpaytma $12$.
\Javob{C}

\Savol{9}
$20d=-60$, $d=-3$, $a_1=70$: \ $a_{101}=-230$, $a_{91}=-200$. \ $\dfrac{230}{200}=1\dfrac3{20}$.
\Javob{C}

\Savol{10}
$t=\log_2x$: \ $t-2$, $t+1$, $t+7$. \ $(t+1)^2=(t-2)(t+7)$ $\Rightarrow$ $2t+1=5t-14$, \ $t=5$, $x=32$.
\Javob{D}

\Savol{11}
\[x\left(x^2+y^2\right)+\frac{xy\left(x^2+y^2\right)}{x-y}=\frac{x^2\left(x^2+y^2\right)}{x-y},\qquad \frac{xy}{x+y}-x=-\frac{x^2}{x+y}.\]
\[-\frac{x^4\left(x^2+y^2\right)}{x^2-y^2}\cdot\frac{x^2-y^2}{x^2\left(x^2+y^2\right)}=-x^2=-9.\]
\Javob{A}

\Savol{12}
$|x-1|>0$ ($x\ne1$), \ $25x-5>0$ $\Rightarrow$ $x>\frac15$: \ $\left(\frac15;1\right)\cup(1;\infty)$.
\Javob{B}

\Savol{13}
$\frac1{x_1}+\frac1{x_2}=\frac{x_1+x_2}{x_1x_2}=6$, \ $\frac1{x_1x_2}=7$.
\[\frac1{x_1^2}+\frac1{x_2^2}=36-14=22,\qquad \frac1{x_1^2x_2^2}=49:\quad t^2-22t+49=0.\]
Koeffitsiyentlar yiğindisi: $1-22+49=28$.
\Javob{C}

\Savol{14}
$u=7^x>0$: \ $7u^3+7\le55u^2-41u$ $\Rightarrow$ $7u^3-55u^2+41u+7\le0$ $\Rightarrow$ $(u-1)(u-7)(7u+1)\le0$.\\
$1\le u\le7$ $\Rightarrow$ $0\le x\le1$: \ $2$ ta butun son.
\Javob{C}

\Savol{15}
$(\sin x-\cos x)^2=2\sin^2\left(x-\frac\pi4\right)=2\sin^2\left(\frac\pi4-x\right)$. \ Natija $2$.
\Javob{A}

\Savol{16}
$L=\log_2x$ ga nisbatan kvadrat: $L^2+(x-1)L+2x-6=0$, \ $D=(x-1)^2-4(2x-6)=(x-5)^2$.\\
$L=-2$ $\Rightarrow$ $x=\frac14$; \ $L=3-x$ $\Rightarrow$ $\log_2x=3-x$, $x=2$ (yagona: çap ösuvçi, öng kamayuvçi).\\
Köpaytma $\frac14\cdot2=\frac12$.
\Javob{A}

\Savol{17}
$(x+1)(x+4)=x^2+5x+4$, \ $(x-2)(x+7)=x^2+5x-14$. \ $t=x^2+5x$: \ $(t+4)(t-14)=19$ $\Rightarrow$ $t^2-10t-75=0$, \ $t=15$ yoki $t=-5$.\\
$x^2+5x-15=0$ va $x^2+5x+5=0$ — ikkalasining ham diskriminanti musbat, har birida yiğindi $-5$. \ Jami $-10$.
\Javob{A}

\Savol{18}
$\cos3x=\cos\left(\frac\pi2-2x\right)$:\\
$3x=\frac\pi2-2x+2\pi k$ $\Rightarrow$ $x=\frac\pi{10}+\frac{2\pi k}5$: \ $-\frac{3\pi}{10},\ \frac\pi{10},\ \frac\pi2,\ \frac{9\pi}{10}$;\\
$3x=-\frac\pi2+2x+2\pi k$ $\Rightarrow$ $x=-\frac\pi2+2\pi k$: \ $-\frac\pi2$. \ Jami $5$ ta.
\Javob{A}

\Savol{19}
$2\le x\le8$ $\Rightarrow$ $t=\frac x2-2\in[-1;2]$, $u=t^2\in[0;4]$: \ $9u-u^2$ — $u\le4{,}5$ da ösuvçi, $u=4$ da $36-16=20$.
\Javob{B}

\Savol{20}
$\dfrac{x^2-xy}{xy}=\dfrac xy-1=2022$.
\Javob{A}

\Savol{21}
\[\frac{2x^2+1}{x^2\left(x^2+1\right)}=\frac1{x^2}+\frac1{x^2+1}\ \Rightarrow\ -\frac1x+\arctg x+C.\]
\Javob{C}

\Savol{22}
$y_1$, $y_2$, $y_4$ — juft; \ $y_3$: $\tg^3x$ toq, $\sqrt{|x|}$ juft — juft ham, toq ham emas. \ $3$ ta.
\Javob{C}

\Savol{23}
$Oy$ ga urinadi: $R=|x_0|=4$.
\Javob{B}

\Savol{24}
$BF=\frac12AB$, $BE=\frac12BC$, \ $\sin\angle FBE=\sin\angle ABC$:
\[S_{BEF}=\frac12\cdot BF\cdot BE\sin\angle FBE=\frac14\cdot\frac12\,AB\cdot BC\sin B=\frac14\cdot60=15.\]
\Javob{A}

\Savol{25}
$y=|(x-1)(x-3)|$: \ $y(2)=1$, $y(3)=0$, $y(4)=3$. \ $\max=3$, $\min=0$, yiğindi $3$.
\Javob{C}

\Savol{26}
Uçburçak tengsizligi: $x+5>3x-5$ $\Rightarrow$ $x<5$; \ $x+3x-5>5$ $\Rightarrow$ $x>2{,}5$. \ Eng katta natural: $4$.
\Javob{A}

\Savol{27}
\fig{\begin{tikzpicture}[scale=0.15,font=\small]
\coordinate (E) at (0,0); \coordinate (M) at (5,0); \coordinate (A) at (15,0); \coordinate (D) at (35,0); \coordinate (F) at (0,5); \coordinate (L) at (5,5); \coordinate (G) at (0,10); \coordinate (N) at (5,10); \coordinate (K) at (15,10); \coordinate (B) at (15,20); \coordinate (C) at (35,20); \coordinate (P) at (11.667,10);
\fill[okgreen!30] (F)--(N)--(C)--cycle;
\draw (E)--(D)--(C)--(B)--(A); \draw (E)--(G)--(N)--(K); \draw (F)--(L); \draw (M)--(N); \draw (F)--(N)--(C)--cycle;
\draw[axs] (N)--(P); \draw[ax] (F)--(35,5);
\node[below] at (E) {$E$}; \node[below] at (M) {$M$}; \node[below] at (A) {$A$}; \node[below] at (D) {$D$}; \node[left] at (F) {$F$}; \node[left] at (G) {$G$}; \node[above] at (N) {$N$}; \node[above] at (B) {$B$}; \node[above right] at (C) {$C$}; \node[aux,below] at (P) {$P$}; \node[right] at (K) {$K$};
\end{tikzpicture}}
Tomonlar: $EF=FL=5$, $FG=5$, $MN=10$, $AK=KB=10$ $\Rightarrow$ $AB=20$.\\
$F$ dan $C$ gaça: gorizontal $5+10+20=35$, vertikal $20-5=15$. $N$ dan $ED$ ga parallel çiziq $FC$ ni $P$ da kesadi; $N$ va $P$ $F$ dan $5$ yuqorida:
\[FP:FC=5:15\ \Rightarrow\ \text{gorizontal böyiça } \frac{35}3,\qquad NP=\frac{35}3-5=\frac{20}3.\]
\[S_{FNC}=\frac12\cdot NP\cdot15=\frac12\cdot\frac{20}3\cdot15=50.\]
\Javob{A}

\Savol{28}
Oltiburçak markazi $O$ bölsa, $S_{BCD}=S_{OBD}$ (rombning yarmi) va hokazo: oltiburçak $BDF$ va unga teng yuzli uçta uçburçak ($ABF$, $BCD$, $DEF$) dan iborat.
\[S_{BDF}=\frac12\cdot36=18.\]
\Javob{B}

\Savol{29}
Kesuvçilar: $CF\cdot CD=CG\cdot CE$ $\Rightarrow$ $6\cdot CD=7\cdot12$, \ $CD=14$.\\
$\triangle CFG\sim\triangle CED$: \ $\dfrac{FG}{ED}=\dfrac{CF}{CE}=\dfrac6{12}$ $\Rightarrow$ $FG=3$.
\Javob{D}

\Savol{30}
$AD\perp$ tekislik: \ $AB=AD\ctg45^\circ=4$, \ $AC=AD\ctg30^\circ=4\sqrt3$.
\[BC^2=16+48-2\cdot4\cdot4\sqrt3\cdot\frac{\sqrt3}2=16,\qquad BC=4.\]
\Javob{D}

\Savol{31}
$A\cap B=\{a,\,b,\,d,\,n\}$: \ $2^4-1=15$.
\Javob{B}

\Savol{32}
Raqamlar $\{6,7,8,9\}$, $3$ ga böliniş qoldiqlari: $6,9\to0$; $7\to1$; $8\to2$. Raqamlar yiğindisi $3$ ga bölinişi kerak.\\
Qoldiqlar: $(0,0,0)$: $2^3=8$; \ $(0,1,2)$ — har xil tartibda: $3!\cdot2=12$; \ $(1,1,1)$: $777$; \ $(2,2,2)$: $888$. \ Jami $8+12+1+1=22$.
\Javob{D}

\newpage
\Savol{33}
\fig{\begin{tikzpicture}[scale=0.75,font=\small]
\draw (0,2) ellipse (1.2 and 0.22); \draw (-1.2,2)--(-1.2,0); \draw (1.2,2)--(1.2,0);
\draw (-1.2,0) arc (180:360:1.2 and 0.22); \draw[dashed] (1.2,0) arc (0:180:1.2 and 0.22);
\draw (-1.2,0) arc (180:360:1.2); \draw[axs] (0,0)--(1.2,0); \draw[axs] (0,0)--(0,-1.2);
\node[left] at (-1.2,2) {$A$}; \node[right] at (1.2,2) {$D$}; \node[left] at (-1.2,0) {$B$}; \node[right] at (1.2,0) {$C$}; \node[aux,above] at (0.6,0) {$R$}; \node[aux,left] at (-1.2,1) {$h$};
\end{tikzpicture}}
Silindr radiusi $R$ (yarim şar radiusi), balandligi $h=AB$.

\textbf{33.} $\dfrac23\pi R^3=\pi R^2h$ $\Rightarrow$ $R=\dfrac32h=30$ cm — \textbf{D}.

\textbf{34.} Ust oçiq: $S=2\pi Rh+2\pi R^2=8\pi+8\pi=16\pi=48$ dm$^2$ — \textbf{A}.

\textbf{35.} $V=\pi R^2h+\dfrac23\pi R^3=8\pi+\dfrac{16\pi}3=\dfrac{40\pi}3=40$ dm$^3=40$ litr — \textbf{C}.
\Javob{33 — D, \ 34 — A, \ 35 — C}

\newpage
\Savol{36}
\[\frac D4=4(3a-2)^2-3\left(a^2+2a\right)=33a^2-54a+16.\]

a) Bitta ildiz: $33a^2-54a+16=0$ — diskriminanti musbat ($54^2-4\cdot33\cdot16=804>0$), Viyet böyiça
\[a_1+a_2=\frac{54}{33}=\frac{18}{11}.\]

b) Yeçim yöq: $33a^2-54a+16<0$ $\Rightarrow$ $a\in\left(\dfrac{27-\sqrt{201}}{33};\ \dfrac{27+\sqrt{201}}{33}\right)\approx(0{,}39;\ 1{,}25)$.\\
Butun qiymat: faqat $a=1$ — $1$ ta.
\Javob{a) $\dfrac{18}{11}$; \quad b) 1}

\newpage
\Savol{37}
Guruhlaymiz: $\sin x+\sin3x+\sin5x=\cos x+\cos3x+\cos5x$.
\[\sin x+\sin5x=2\sin3x\cos2x,\qquad \cos x+\cos5x=2\cos3x\cos2x:\]
\[\sin3x(1+2\cos2x)=\cos3x(1+2\cos2x)\ \Rightarrow\ (1+2\cos2x)(\sin3x-\cos3x)=0.\]
$\cos2x=-\frac12$ $\Rightarrow$ $x=\pm\frac\pi3+\pi k$: \ $-\frac\pi3,\ \frac\pi3,\ \frac{2\pi}3$.\\
$\tg3x=1$ $\Rightarrow$ $x=\frac\pi{12}+\frac{\pi k}3$: \ $-\frac\pi4,\ \frac\pi{12},\ \frac{5\pi}{12},\ \frac{3\pi}4$.

a) $3+4=7$ ta.

b) $\left(-\frac\pi3+\frac\pi3+\frac{2\pi}3\right)+\left(-\frac{3\pi}{12}+\frac\pi{12}+\frac{5\pi}{12}+\frac{9\pi}{12}\right)=\frac{2\pi}3+\pi=\frac{5\pi}3$.
\Javob{a) 7; \quad b) $\dfrac{5\pi}3$}

\newpage
\Savol{38}
$f'(x)=3x^2+a$. Urinmalar parallel: $f'(a)=f'(b)$ $\Rightarrow$ $a^2=b^2$, $a\ne b$ $\Rightarrow$ $b=-a$.

a) $f(1)=1+a+b=1$.

b) $b=2$ $\Rightarrow$ $a=-2$: \ $f'(1)=3-2=1$.
\Javob{a) 1; \quad b) 1}

\newpage
\Savol{39}
\fig{\begin{tikzpicture}[x=0.65cm,y=0.45cm,font=\small]
\fill[okgreen!30] (2,0)--(2,2.4) plot[smooth,tension=0.8] coordinates {(2,2.4) (2.15,3.1) (3,3.5) (4.2,4.0) (5,5)}--(5,0)--cycle;
\fill[red!15] (0,2.4)--(2,2.4) plot[smooth,tension=0.8] coordinates {(2,2.4) (2.15,3.1) (3,3.5) (4.2,4.0) (5,5)}--(0,5)--cycle;
\draw[-{Stealth[length=4pt]}] (-0.3,0)--(7,0) node[right] {$x$}; \draw[-{Stealth[length=4pt]}] (0,-0.3)--(0,6.6) node[left] {$y$};
\draw[thick] plot[smooth,tension=0.8] coordinates {(2,2.4) (2.15,3.1) (3,3.5) (4.2,4.0) (5,5) (5.5,6.2)};
\draw[dashed] (2,0)--(2,2.4)--(0,2.4); \draw[dashed] (5,0)--(5,5)--(0,5);
\node[below] at (2,0) {$2$}; \node[below] at (5,0) {$5$}; \node[left] at (0,2.4) {$a$}; \node[left] at (0,5) {$b$};
\end{tikzpicture}}
$\displaystyle\int_2^5f\,f'\,dx=\frac{f^2(5)-f^2(2)}2=\frac{b^2-a^2}2=13{,}5$ $\Rightarrow$ $(b-a)(b+a)=27$, \ $b-a=3$.\\
$a=3$, $b=6$.

a) $\dfrac{a-b}2=-\dfrac32$.

b) Rasmda: $\int_2^5f\,dx$ — yaşil soha, $\int_a^bf^{-1}dx$ — puşti soha (egri çiziqning $Oy$ tomoni). Ikkalasi birga katta töğri törtburçakdan kiçigini ayirganiga teng:
\[5b-2a=30-6=24.\]
\Javob{a) $-\dfrac32$; \quad b) 24}

\newpage
\Savol{40}
\fig{\begin{tikzpicture}[scale=0.95,font=\small]
\draw[-{Stealth[length=4pt]}] (-1.9,0)--(2,0) node[right] {$x$}; \draw[-{Stealth[length=4pt]}] (0,-0.2)--(0,4.6) node[left] {$y$};
\draw[thick,domain=-2.05:2.05,samples=60] plot (\x,{\x*\x});
\draw (0,0.61) circle (0.6); \draw (0,2.81) circle (1.6);
\coordinate (A) at (0.332,0.11); \coordinate (B) at (1.52,2.31); \coordinate (O1) at (0,0.61); \coordinate (O2) at (0,2.81); \coordinate (C) at (0,1.21);
\draw[axs] (O1)--(A); \draw[axs] (O2)--(B); \draw[ax] (0,0.11)--(A); \draw[ax] (0,2.31)--(B);
\fill (A) circle (1.3pt); \fill (B) circle (1.3pt); \fill (O1) circle (1.3pt); \fill (O2) circle (1.3pt); \fill (C) circle (1.3pt);
\node[right] at (A) {$A$}; \node[right] at (B) {$B$}; \node[left] at (O1) {$O_1$}; \node[left] at (O2) {$O_2$}; \node[above left] at (C) {$C$};
\end{tikzpicture}}
Uriniş nuqtasi $A\left(x_1;x_1^2\right)$. Radius $O_1A$ urinmaga perpendikulyar (normal). Urinma oğmasi $2x_1$, normal oğmasi $-\frac1{2x_1}$:
\[O_1(0;\,y_0):\quad \frac{x_1^2-y_0}{x_1-0}=-\frac1{2x_1}\ \Rightarrow\ y_0=x_1^2+\frac12=y_1+\frac12.\]

a) $O_1$ va $(0;y_1)$ orasidagi masofa $\dfrac12$.

b) $r^2=x_1^2+\left(\frac12\right)^2=y_1+\frac14$ $\Rightarrow$ markaz ordinatasi $y_0=y_1+\frac12=r^2+\frac14$. Xuddi şunday $O_2$ ordinatasi $R^2+\frac14$.\\
Aylanalar taşqi urinadi: $O_1O_2=R+r$:
\[R^2-r^2=R+r\ \Rightarrow\ R-r=1,\qquad O_2B=R=2024.\]
\Javob{a) $\dfrac12$; \quad b) 2024}

\newpage
\Savol{41}
\fig{\begin{tikzpicture}[scale=0.62,font=\small]
\coordinate (A) at (0,0); \coordinate (D) at (7,0); \coordinate (B) at (3,4); \coordinate (C) at (10,4);
\coordinate (H) at (4.4,4); \coordinate (E) at (5.6,0); \coordinate (G) at (7,4); \coordinate (L) at (3,0); \coordinate (O) at (5,2);
\fill[okgreen!30] (H)--(G)--(O)--cycle; \fill[okgreen!30] (L)--(E)--(O)--cycle;
\draw (A)--(B)--(C)--(D)--cycle; \draw (H)--(E); \draw (L)--(G); \draw[ax] (A)--(C); \draw[ax] (B)--(D);
\node[below left] at (A) {$A$}; \node[above] at (B) {$B$}; \node[above right] at (C) {$C$}; \node[below right] at (D) {$D$};
\node[above] at (H) {$H$}; \node[above] at (G) {$G$}; \node[below] at (L) {$L$}; \node[below] at (E) {$E$}; \node[left] at (4.85,2.1) {$O$};
\end{tikzpicture}}
$BH=ED$ va $BH\parallel ED$ $\Rightarrow$ $HE$ parallelogramm markazidan ötadi; xuddi şunday $AL=CG$ $\Rightarrow$ $LG$ ham markazdan ötadi. Demak $O$ — markaz, $\triangle HOG=\triangle EOL$.\\
$BC=a$, balandlik $h$, $S=ah$. \ $HG=a-\frac a5-\frac{3a}7=\frac{13}{35}a$; $O$ dan $BC$ gaça $\frac h2$.

a) \[2\cdot\frac12\cdot\frac{13}{35}a\cdot\frac h2=\frac{13}{70}S.\]

b) $HE$ parallelogrammni teng ikkiga böladi: $S_{ABHE}=\frac S2$. \ $S_{ABHOL}=S_{ABHE}-S_{LOE}$:
\[\frac12-\frac{13}{140}=\frac{57}{140}.\]
\Javob{a) $\dfrac{13}{70}$; \quad b) $\dfrac{57}{140}$}

\newpage
\Savol{42}
\fig{\begin{tikzpicture}[scale=0.5,font=\small]
\coordinate (B) at (0,0); \coordinate (C) at (8,0); \coordinate (A) at (2.5,4.33); \coordinate (I) at (3,1.732);
\coordinate (K) at (3,0); \coordinate (M) at (1.5,2.598); \coordinate (N) at (4.071,3.093);
\draw (A)--(B)--(C)--cycle; \draw (I) circle (1.732); \draw[axs] (M)--(N);
\fill (K) circle (2.5pt); \fill (M) circle (2.5pt); \fill (N) circle (2.5pt);
\node[above] at (A) {$A$}; \node[below left] at (B) {$B$}; \node[below right] at (C) {$C$}; \node[below] at (K) {$K$}; \node[left] at (M) {$M$}; \node[above right] at (N) {$N$};
\node[below] at (1.5,0) {$3$}; \node[below] at (5.5,0) {$5$}; \node[aux,left] at (2,3.5) {$x$}; \node[aux,above] at (3.3,3.8) {$x$};
\end{tikzpicture}}
Urinma kesmalari: $BM=BK=3$, $CN=CK=5$, $AM=AN=x$. \ $AB=x+3$, $AC=x+5$, $BC=8$.\\
Kosinuslar teoremasi ($\angle B=60^\circ$):
\[(x+5)^2=(x+3)^2+64-8(x+3)\ \Rightarrow\ 12x=24,\ x=2.\]
$AB=5$, $AC=7$.

a) $S=\frac12\cdot5\cdot8\cdot\frac{\sqrt3}2=10\sqrt3$.

b) $\cos A=\dfrac{25+49-64}{2\cdot5\cdot7}=\dfrac17$. \ $\triangle AMN$: $AM=AN=2$:
\[MN^2=4+4-2\cdot2\cdot2\cdot\frac17=\frac{48}7,\qquad MN=\sqrt{\frac{48}7}.\]
\Javob{a) $10\sqrt3$; \quad b) $\sqrt{\dfrac{48}7}$}

\newpage
\Savol{43}
\fig{\begin{tikzpicture}[scale=2.1,font=\small]
\coordinate (B) at (0,0); \coordinate (C) at (1,0); \coordinate (D) at (1.309,0.951); \coordinate (E) at (0.5,1.539); \coordinate (A) at (-0.1045,0.9945);
\draw (A)--(B)--(C)--(D)--(E)--cycle; \draw[axs] (B)--(E); \draw[ax] (B)--(D);
\node[left] at (A) {$A$}; \node[below left] at (B) {$B$}; \node[below right] at (C) {$C$}; \node[right] at (D) {$D$}; \node[above] at (E) {$E$};
\end{tikzpicture}}
$BC=CD=DE=s$, $\angle C=\angle D=108^\circ$ — muntazam beşburçakning $4$ ta ketma-ket uçi kabi:
\[\triangle BCD:\ \angle CBD=\angle CDB=36^\circ;\qquad \triangle CDE:\ \angle DCE=\angle DEC=36^\circ.\]
$BCDE$ — teng yonli trapetsiya ($BE\parallel CD$), \ $\angle CBE=\angle DEB=72^\circ$, \ $BE=BD=2s\cos36^\circ$.\\
$\angle ABE=96^\circ-72^\circ=24^\circ$. \ $\triangle ABE$: $AB=s$, $BE=2s\cos36^\circ$, $\angle AEB=y$, $\angle BAE=156^\circ-y$. Sinuslar teoremasi:
\[\frac{\sin(156^\circ-y)}{\sin y}=\frac{BE}{AB}=2\cos36^\circ.\]
$y=30^\circ$ da: $\dfrac{\sin126^\circ}{\sin30^\circ}=2\sin54^\circ=2\cos36^\circ$ — bajariladi (çap tomon $y$ ga nisbatan kamayuvçi, yeçim yagona).

a) $\angle A=\angle BAE=156^\circ-30^\circ=126^\circ$.

b) $\angle E=\angle AEB+\angle BED=30^\circ+72^\circ=102^\circ$.
\Javob{a) $126^\circ$; \quad b) $102^\circ$}

\newpage
\Savol{44}
$V_{\text{um}}=\dfrac{V_{\text{ş}}}{12}=\dfrac{V_k}{16}$ $\Rightarrow$ $V_{\text{ş}}=\dfrac34V_k$; \qquad $\dfrac{V_t}{10}=\dfrac{V_k}{20}$ $\Rightarrow$ $V_t=\dfrac12V_k$.

a) $\dfrac{V_{\text{ş}}}{V_t}=\dfrac{3/4}{1/2}=\dfrac32$.

b) Muntazam tetraedr: $V_t=\dfrac{a^3}{6\sqrt2}=\dfrac{432\sqrt2}{6\sqrt2}=72$.\\
$V_k=2V_t=144$.
\Javob{a) $\dfrac32$; \quad b) 144}

\newpage
\Savol{45}
\fig{\begin{tikzpicture}[scale=0.85,font=\small]
\fill[okgreen!30] (-1.1,2.6)--(0,0)--(1.1,2.6) arc (0:-180:1.1 and 0.22);
\draw (0,2.6) ellipse (1.1 and 0.22); \draw (-1.1,2.6)--(0,0)--(1.1,2.6); \draw[axs] (0,0)--(0,2.6)--(1.1,2.6);
\draw (0,2.4)--(0.2,2.4)--(0.2,2.6);
\node[aux,left] at (0,1.2) {$H$}; \node[right] at (0.6,1.3) {$R$}; \node[aux,above] at (0.55,2.6) {$r$};
\end{tikzpicture}}
Konus yasovçisi $R$; asos aylanasi uzunligi qolgan yoy uzunligiga teng:
\[2\pi r=\frac{360^\circ-\alpha}{360^\circ}\cdot2\pi R\ \Rightarrow\ r=\frac{360^\circ-\alpha}{360^\circ}R.\]

a) $\alpha=60^\circ$: $r=\frac56\cdot12=10$, \ $H=\sqrt{144-100}=2\sqrt{11}$:
\[V=\frac13\pi r^2H=\frac13\cdot3\cdot100\cdot2\sqrt{11}=200\sqrt{11}.\]

b) $V(r)=\dfrac\pi3r^2\sqrt{R^2-r^2}$. \ $V^2\sim r^4\left(R^2-r^2\right)$; \ $\left(r^4R^2-r^6\right)'=4r^3R^2-6r^5=0$ $\Rightarrow$ $r=R\sqrt{\frac23}=\dfrac{\sqrt6}3R$.
\[360^\circ-\alpha=360^\circ\cdot\frac{\sqrt6}3\approx120^\circ\cdot2{,}45=294^\circ,\qquad \alpha\approx66^\circ.\]
\Javob{a) $200\sqrt{11}$; \quad b) $66^\circ$}

\ShowJavoblar
\end{document}
